\chapter{Das Erbtheil der Heiligen im Licht}
\quellseite{9}

„Würde es nicht jeder für eine große Wohlthat achten, wenn er in das Erbtheil des reichsten und mächtigsten Königs eingesetzt würde? Nun ist es aber eine viel größere Wohlthat, daß wir Theil und Gemeinschaft haben an dem Erbe der heiligen Patriarchen, Propheten und Apostel. Es hat Leute gegeben, die frei öffentlich gesagt haben: gern wollten sie bis zum jüngsten Tag in der Hölle sitzen, wenn sie nur gewiß wären, daß sie dann ewig selig werden könnten. Ach, die armen Leute hätten ja aus der Predigt des Evangelii den Trost haben können, daß wir der Seligkeit und des Erbtheils der Heiligen durch unsern Herrn Christum gewiß sein dürfen. Ebenso sind Leute gewesen, die an ihrem letzten Ende all ihr Hab und Gut willig hingegeben hätten, wenn sie des Himmelreichs und der Seligkeit hätten gewiß sein können. Aber im Evangelio ist uns ja die Thür zum Himmelreich weit genug aufgethan und gesagt, daß wir ohne all unser Verdienst durch Christum zur Seligkeit kommen können. Darum können wir für solche große Wohlthat nimmermehr genug danken.“

(Joh. Brenz.)
